\section{Discussion}\label{discussion}

Long-read transcriptome sequencing has transformed transcriptome analysis by enabling direct observation of full-length isoforms rather than inference from short-read fragments. Recent gains in throughput and accuracy have extended these capabilities to single-cell and single-nucleus settings, creating new opportunities to study isoform regulation in heterogeneous tissues and disease. At the same time, these advances have exposed persistent computational challenges, including partial transcript coverage, ambiguous read-to-isoform compatibility, and dependence on incomplete reference annotations. Addressing these challenges requires methods that integrate accurate isoform discovery with principled quantification and that generalize across sequencing platforms and experimental contexts.

Here we present LRAA as a unified framework for isoform identification and quantification from long-read RNA sequencing data spanning bulk, single-cell, and single-nucleus applications. LRAA integrates splice-graph--based structural modeling with an expectation--maximization framework for abundance estimation, enabling probabilistic resolution of ambiguous reads while retaining sensitivity to novel isoforms. Across a broad range of benchmarking datasets, LRAA demonstrated robust performance across execution modes, sequencing technologies, and biological settings. In evaluations that included leading methods assessed by the LRGASP Consortium, such as IsoQuant, Bambu and StringTie, LRAA ranked among the strongest-performing approaches in quantification-only, reference-guided, and de novo analyses. In cases where LRAA did not achieve the highest overall accuracy, its performance closely matched that of the top-ranked method.

Rigorous benchmarking of long-read isoform analysis methods remains challenging because no single dataset simultaneously provides complete ground truth and biological realism. Simulated datasets offer unambiguous isoform-of-origin information and are valuable for controlled evaluation of algorithmic behavior, including sensitivity to sequencing error. However, simulations necessarily lack native transcriptional processing and degradation. Synthetic spike-in controls such as SIRVs incorporate real sequencing characteristics and defined expression ratios, but remain artificial constructs that do not fully reflect cellular transcript processing. As a result, method performance on simulations or spike-ins does not always generalize to biologically complex datasets.

To address this gap, we leveraged the MORF library as a biologically informed benchmarking substrate. MORFs consist of thousands of human isoforms expressed in living cells, each tagged with an isoform-specific barcode that enables unambiguous assignment of every read, including partial reads, to its isoform of origin. This design combines the ground-truth clarity of simulations with the biological context of native transcription. MORF benchmarking revealed certain tools to have notable differences in relative method performance compared with benchmarks on SIRVs, while others, such as LRAA, maintained consistently strong performance across benchmarking conditions. These results underscore the value of biologically inspired benchmarks for realistic evaluation of long-read isoform analysis tools.

Beyond controlled benchmarks, we evaluated LRAA on whole-transcriptome datasets from multiple human cell lines and reference RNA samples sequenced using both PacBio and ONT technologies. In this setting, where absolute ground truth is unavailable, LRAA showed strong performance for isoform identification and produced expression profiles that closely matched those of methods previously shown to perform well under benchmarks with defined truth sets. These results indicate that LRAA generalizes effectively to real-world transcriptome analyses and performs robustly across sequencing platforms without requiring platform-specific workflows.

A major advance of this work is the extension of isoform-resolved long-read analysis to single-cell and single-nucleus transcriptomics. Single-cell long-read data are sparse, and biologically relevant isoforms may be confined to specific cell populations and obscured in pseudobulk analyses. By integrating isoform discovery with cell clustering, LRAA enables recovery of isoforms that are low abundance at the sample level but prominent within particular cell types. In PBMCs, this approach revealed cell-type-specific isoform usage driven by both alternative splicing and alternative transcript start or end site selection.

Single-nucleus transcriptomics presents additional challenges, as nuclear RNA is enriched for unspliced and partially processed transcripts. Despite these constraints, LRAA identified biologically meaningful isoform usage in single-nucleus data from FTD frontal cortex tissue. In particular, we detected the pathogenic cryptically spliced isoform of \emph{STMN2}, a hallmark of TDP-43 proteinopathy, with marked differences in prevalence across neuronal subpopulations. These observations demonstrate that long-read single-nucleus sequencing, coupled with appropriate computational modeling, can resolve disease-relevant isoform dysregulation in tissues that are poorly represented in conventional single-cell approaches. \tofix{Transcript termini need more care in nuclei. Reads from the FTD nuclei, and neuronal nuclei in particular, rarely reached the 5\textquotesingle{} end of their transcripts, and the apparent TSS switches between neurons and other cells were dominated by downstream sites without CAGE support, whereas PolyA sites were as reliable as in PBMCs. Template-switching library preparation leaves an extra, untemplated G at the 5\textquotesingle{} end of molecules copied from capped RNA \citep{wulf2019, cumbie2015}, and in the Kinnex single-cell data this signature separated reads starting at a TSS from truncated reads and marked the FTD downstream sites as truncation. The signature is not specific to Kinnex single-cell libraries: in the bulk PacBio and ONT cDNA libraries benchmarked here, 74\% and 28--29\% of reads starting at annotated TSSs carried one G more than the library's baseline, against 6--7\% and 3\% of reads starting inside transcript bodies, whereas the 5\textquotesingle{} ends of the uncapped SIRV spike-ins in the same PacBio libraries carried it no more often than their transcript bodies (3\%). The baseline itself differed between protocols (no G for single-cell Kinnex, one for bulk PacBio, and three for ONT reads retaining the strand-switching primer), and it is lost where 5\textquotesingle{} clips are trimmed before alignment. A future version of LRAA could use this signature, calibrated per library against reads starting inside transcript bodies, when calling TSS sites and pruning 5\textquotesingle{} degradation products, as Ouro-Seq does for nanopore single-cell data \citep{an2025}.}{site-first restructure, 2026-10-27: cap-signature discussion; cross-dataset numbers from \_\_TSS\_5prime\_Completeness/cap\_signature\_other\_datasets (BT474/K562 bulk PacBio, MCF7/A549 SG-NEx PCR-cDNA, SIRVs; MORF clips trimmed upstream)}

Our identification of novel proteoforms, protein isoforms arising from unannotated splice junctions or intron retention events, revealed that \textasciitilde15\% of the PBMC proteome and \textasciitilde6\% of the FTD frontal cortex tissue proteome may consist of sequences not present in current reference annotations. These novel proteoforms represent complete ORFs with canonical start and stop codons, suggesting translational competence, though translatomics and proteomics validation will be essential to confirm bona fide expressed proteins and to establish their functional relevance in health and disease.

Several limitations remain. Although LRAA models alternative transcription start and end sites, benchmarking here focused primarily on splice-pattern accuracy, reflecting ongoing challenges in reliably assessing transcript termini in the presence of RNA degradation and protocol-specific biases. Single-exon isoforms remain difficult to evaluate in de novo discovery settings, necessitating conservative filtering to limit false positives. Continued improvements in sequencing protocols, library preparation, and benchmarking standards will further advance isoform-resolved transcriptomics.

As long-read sequencing is increasingly integrated with single-cell, single-nucleus, and emerging spatial technologies, the ability to accurately resolve isoform structure and usage at scale will become central to transcriptome analysis. By combining robust isoform discovery, principled quantification, and applicability across experimental contexts, LRAA represents a foundational framework for dissecting isoform regulation across development and disease.
